\subsection{\texorpdfstring{COMMUTANT OF \(\mathfrak{h}\) IN THE ENVELOPING ALGEBRA OF \(\mathfrak{g}\)}{COMMUTANT OF \textbackslash mathfrak\{h\} IN THE ENVELOPING ALGEBRA OF \textbackslash mathfrak\{g\}}}

Let U be the enveloping algebra of g, \(V \subset U\) the enveloping algebra of h. The algebra V can be identified with the symmetric algebra \(\mathbf{S}(\mathfrak{h})\) of h, and also with the algebra of polynomial functions on \(h^{*}\) . Denote by \(\alpha_{1},\ldots,\alpha_{n}\) the pairwise distinct positive roots. Let \((H_{1},\ldots,H_{l})\) be a basis of h. By the Poincaré-Birkhoff-Witt theorem, the elements

\[
u ((q _ {i}), (m _ {i}), (p _ {i})) = X _ {- \alpha_ {1}} ^ {q _ {1}} \dots X _ {- \alpha_ {n}} ^ {q _ {n}} H _ {1} ^ {m _ {1}} \dots H _ {l} ^ {m _ {l}} X _ {\alpha_ {1}} ^ {p _ {1}} \dots X _ {\alpha_ {n}} ^ {p _ {n}}
\]

\((q_{i},m_{i},p_{i})\) integers \(\geq 0\) ) form a basis of the vector space U. For all \(h\in \mathfrak{h}\) , we have

\[
[ h, u ((q _ {i}), (m _ {i}), (p _ {i})) ] = ((p _ {1} - q _ {1}) \alpha_ {1} + \dots + (p _ {n} - q _ {n}) \alpha_ {n}) (h) u ((q _ {i}), (m _ {i}), (p _ {i})).\tag{3}
\]

The vector space U is a g-module (hence also an h-module) under the adjoint representation. If \(\lambda \in h^{*}\) , the subspaces \(U^{\lambda}\) and \(U_{\lambda}\) are defined (Chap. VII, §1, no. 3); formula (3) shows that \(U^{\lambda} = U_{\lambda}\) and that \(U = \bigoplus_{\lambda \in Q} U^{\lambda}\) (where Q is the group of radical weights of R). In particular, \(U^{0}\) is the commutant of h, or of V, in U.

Lemma 3. Put \(\mathrm{L} = (\mathfrak{n}_{-}\mathrm{U})\cap \mathrm{U}^{0}\)

\begin{enumerate}
\def\labelenumi{(\roman{enumi})}
\item
  We have \(\mathrm{L} = (\mathrm{U}\mathfrak{n}_{+})\cap \mathrm{U}^{0}\) , and \(\mathrm{L}\) is a two-sided ideal of \(\mathrm{U}^0\) .
\item
  We have \(U^{0} = V \oplus L\) .
\end{enumerate}

It is clear that \(\mathfrak{n}_{-}\mathrm{U}\) (resp. \(\mathrm{U}\mathfrak{n}_{+}\) ) is the set of linear combinations of the elements \(u((q_{i}),(m_{i}),(p_{i}))\) such that \(\sum q_{i}>0\) (resp. \(\sum p_{i}>0\) ). On the other hand

\[
u \left(\left(q _ {i}\right), \left(m _ {i}\right), \left(p _ {i}\right)\right) \in \mathrm{U} ^ {0} \Longleftrightarrow p _ {1} \alpha_ {1} + \dots + p _ {n} \alpha_ {n} = q _ {1} \alpha_ {1} + \dots + q _ {n} \alpha_ {n}.
\]

This implies that \((\mathfrak{n}_{-}\mathrm{U})\cap\mathrm{U}^{0}=(\mathrm{U}\mathfrak{n}_{+})\cap\mathrm{U}^{0}\) . Finally, \((\mathfrak{n}_{-}\mathrm{U})\cap\mathrm{U}^{0}\) (resp. \((\mathrm{U}\mathfrak{n}_{+})\cap\mathrm{U}^{0}\) ) is a right (resp. left) ideal of \(U^{0}\) , hence (i). Further, an element \(u((q_{i}),(m_{i}),(p_{i}))\) that is in \(U^{0}\) belongs to V (resp. to L) if and only if \(p_{1}=\cdots=p_{n}=q_{1}=\cdots=q_{n}=0\) (resp. \(p_{1}+\cdots+p_{n}+q_{1}+\cdots+q_{n}>0\) ), hence (ii). Q.E.D.

In view of Lemma 3, the projection of \(U^{0}\) onto V with kernel L is a homomorphism of algebras. It is called the Harish-Chandra homomorphism from \(U^{0}\) to V (relative to B). Recall that V can be identified with the algebra of polynomial functions on \(h^{*}\) .

PROPOSITION 7. Let \(\lambda \in \mathfrak{h}^*\) , E a \(\mathfrak{g}\) -module generated by a primitive element of weight \(\lambda\) , \(\chi\) the central character of E, and \(\varphi\) the Harish-Chandra homomorphism from \(\mathrm{U}^0\) to V. Then, \(\chi(z) = (\varphi(z))(\lambda)\) for all \(z\) in the centre of U.

Let v be a primitive element of E of weight \(\lambda\) , and z an element of the centre of U. There exist \(u_{1},\ldots,u_{p}\in U\) and \(n_{1},\ldots,n_{p}\in\mathfrak{n}_{+}\) such that \(z=\varphi(z)+u_{1}n_{1}+\cdots+u_{p}n_{p}\) . Then

\[
\chi (z) v = z v = \varphi (z) v + u _ {1} n _ {1} v + \dots + u _ {p} n _ {p} v = \varphi (z) v = (\varphi (z)) (\lambda) v.
\]

COROLLARY. Let \(\langle\cdot,\cdot\rangle\) be a non-degenerate invariant symmetric bilinear form on g, C the Casimir element associated to \(\langle\cdot,\cdot\rangle\) . Denote also by \(\langle\cdot,\cdot\rangle\) the inverse form on \(h^{*}\) of the restriction of \(\langle\cdot,\cdot\rangle\) to h (§2, no. 3, Prop. 5). Then \(\chi(\mathrm{C})=\langle\lambda,\lambda+2\rho\rangle\) , where \(\rho=\frac{1}{2}\sum_{\alpha\in R_{+}}\alpha\) .

We recall the notations of §2, no. 3, Prop. 6. We have

\[
\begin{array}{r} \mathrm{C} = \sum_ {\alpha \in \mathrm{R} _ {-}} \frac {1}{\langle X _ {\alpha} , X _ {- \alpha} \rangle} X _ {\alpha} X _ {- \alpha} + \sum_ {\alpha \in \mathrm{R} _ {+}} \frac {1}{\langle X _ {\alpha} , X _ {- \alpha} \rangle} X _ {- \alpha} X _ {\alpha} \\ + \sum_ {\alpha \in \mathrm{R} _ {+}} \frac {1}{\langle X _ {\alpha} , X _ {- \alpha} \rangle} [ X _ {\alpha}, X _ {- \alpha} ] + \sum_ {i \in I} H _ {i} H _ {i} ^ {\prime} \end{array}
\]

so

\[
\varphi (\mathrm{C}) = \sum_ {\alpha \in \mathrm{R} _ {+}} \frac {1}{\langle X _ {\alpha} , X _ {- \alpha} \rangle} [ X _ {\alpha}, X _ {- \alpha} ] + \sum_ {i \in I} H _ {i} H _ {i} ^ {\prime}.
\]

By Prop. 7,

\[
\chi (\mathrm{C}) = \sum_ {\alpha \in \mathrm{R} _ {+}} \frac {1}{\langle X _ {\alpha} , X _ {- \alpha} \rangle} \lambda ([ X _ {\alpha}, X _ {- \alpha} ]) + \sum_ {i \in I} \lambda (H _ {i}) \lambda (H _ {i} ^ {\prime}).
\]

Let \(h_{\lambda}\) be the element of \(\mathfrak{h}\) such that \(\langle h_{\lambda}, h \rangle = \lambda(h)\) for all \(h \in \mathfrak{h}\) . By §2, no. 2, Prop. 1,

\[
\lambda \left(\frac {1}{\langle X _ {\alpha} , X _ {- \alpha} \rangle} [ X _ {\alpha}, X _ {- \alpha} ]\right) = \left\langle h _ {\lambda}, \frac {1}{\langle X _ {\alpha} , X _ {- \alpha} \rangle} [ X _ {\alpha}, X _ {- \alpha} ] \right\rangle = \alpha (h _ {\lambda}) = \langle \lambda , \alpha \rangle .
\]

Hence

\[
\chi (\mathrm{C}) = \left(\sum_ {\alpha \in \mathrm{R} _ {+}} \langle \lambda , \alpha \rangle\right) + \langle \lambda , \lambda \rangle = \langle \lambda , \lambda + 2 \rho \rangle .
\]

